% Generated draft: requires mathematical and source review.
\documentclass[11pt]{article}
\usepackage[margin=25mm]{geometry}
\usepackage{fontspec}
\setmainfont{Latin Modern Roman}
\usepackage{amsmath,amssymb,mathtools,unicode-math}
\setmathfont{Latin Modern Math}
\usepackage{enumitem,needspace,xurl,hyperref}
\setlength{\parindent}{0pt}
\setlength{\parskip}{5pt}
\setlength{\emergencystretch}{3em}
\newcommand{\N}{\mathbb{N}}
\newcommand{\Z}{\mathbb{Z}}
\newcommand{\Q}{\mathbb{Q}}
\newcommand{\R}{\mathbb{R}}
\newcommand{\C}{\mathbb{C}}
\newcommand{\F}{\mathbb{F}}
\newcommand{\A}{\mathbb{A}}
\newcommand{\PP}{\mathbb{P}}
\newcommand{\E}{\mathbb{E}}
\newcommand{\eps}{\varepsilon}
\newcommand{\dd}{\,\mathrm{d}}
\begin{document}
\section*{UnsolvedMath ID: 30001875}
\textbf{Source problem code:} catalogue entry
\subsection{Definitions and mathematical statement}
For a positive integer \(k\), let \(S_{2k}(\operatorname{SL}_2(\Z))\) denote the space of holomorphic cusp forms of weight \(2k\) for \(\operatorname{SL}_2(\Z)\). A Hecke eigenform is a nonzero form that is a simultaneous eigenvector for the Hecke operators. The supplied question takes a normalized Hecke eigenform \(f\) in this space and denotes its Saito--Kurokawa lift by \(F_f\).

The source asks whether one can prove unconditionally that
\[
N(F_f)\sim 2\qquad\text{as }k\to\infty.
\]
Here the asserted asymptotic means that \(N(F_f)/2\) tends to \(1\) along the intended family of forms and weights.

The supplied record does not define the quantity \(N(F_f)\), its measures or normalization, the normalization of \(f\) and \(F_f\), or the weights for which the intended lift is defined. It also does not specify how \(f\) varies with \(k\), or whether the assertion is uniform over all eligible eigenforms. Consequently, the displayed question preserves the supplied target but is not a fully specified mathematical proposition. These missing definitions and quantifiers must be recovered from the source before a precise formulation can be finalized.
\subsection{Short English statement}
The question asks for an unconditional proof that the source's quantity \(N(F_f)\), attached to a Saito--Kurokawa lift of a normalized Hecke eigenform of increasing weight, tends to \(2\). The quantity, normalization conventions, eligible weights, and quantification over eigenforms are not specified in the supplied record and require clarification.
\subsection{Sources}
\begin{itemize}
\item Supplied local registry record: permanent UnsolvedMath ID \(30001875\), title ``Asymptotic Norms of Saito--Kurokawa Lifts,'' source problem code \texttt{catalogue entry}, in the Oberwolfach Reports Open Problems collection. No individual proposer is supplied; the recorded proposed year is 2011.
\item Supplied source citation: \emph{The Analytic Theory of Automorphic Forms} (2012), Oberwolfach Reports, pp.\ 2397--2458, DOI \texttt{10.4171/owr/2011/43}. Original OWR report: \url{https://doi.org/10.4171/owr/2011/43}.
\item Supplied literature reference, labeled ``second moments in the weight aspect'': \url{https://msp.org/ant/2018/12-1/ant-v12-n1-p.pdf}. No authors, article title, or precise supporting result are supplied.
\item Literature assessment checked at \texttt{2026-08-21}: the supplied curation labels the problem open and states that the symmetric-square second-moment problem in the weight aspect was open in the source and remains unresolved in general despite bounds and special averages. This is supplied dated curation, not an independently verified literature assessment. The record expressly requests independent verification of time-sensitive status and cited claims.
\item The supplied statement status is \texttt{corrected\_verified}; its verification note reports comparison with the official Oberwolfach report and normalization of notation, context, and question form without changing the source-level claim. This is a supplied verification claim, not independent verification performed for this draft.
\end{itemize}
\end{document}
